
\fancyhf{}
% \lhead{\begin{tabular}{|p{8.25cm}} \hline {\large \bf SNO+ Laser Procedures} \\ \\ \\ \\ \end{tabular}}
\lhead{\begin{tabular}{|p{8cm}|p{3cm}|p{3cm}|}  
	\hline 
		{\large \bf SNO+ Cleaning Module Procedures}  
										& Document No:  & Rev. No: 01\\ 
								\cline{2-3} & Date: 2019-07-07 & 
								Page \thepage \\  \hline \end{tabular}}


\begin{tabular}{||l|l|l||}
\hline\hline
& \multicolumn{2}{p{8cm}||}{\bf SNO+ Cleaning Module Procedures} \\
\includegraphics[width=6cm]{figures/SNOplus_logo.png} & \multicolumn{2}{p{8cm}||}{} \\
\hline
\multicolumn{2}{||p{8.5cm}|}{Document Number:} & Revision Number: 01\\
\hline
\multicolumn{3}{||l||}{Document Owner: SNO+ Calibration Post-Doc} \\
\hline
\multicolumn{3}{||l||}{Reviewer:}\\
\hline
Name: & Signature & Date \\
\hline
\multicolumn{3}{||l||}{Authorizer:}\\
\hline
Name: & Signature & Date \\
\hline\hline
\end{tabular}
\thispagestyle{empty}



\section{Purpose}

The cleaning module is a "hands off" means of cleaning the calibration sources and evaluating the cleanliness of the source. However, to connect and disconnect the source from the umbilical it is impractical to make use of the source connector in a glove box, so it is expected that the connection will be made with the source and umbilical exposed to the DCR air. Special steps must be taken to mitigate the effect of this exposure on the source, URM and experiment at large.

\subsection{Definitions}
\begin{itemize}
\item {\bf Cleaning Module} An apparatus that can be attached directly to the source tube gate valve that sprays a flow of LAB directly onto a source to facilitate cleaning and a drain from which samples of LAB can be drawn for testing.
\item {\bf Storage Module} An air tight apparatus in which a source may be stored between deployments when disconnected from the URM.
\item {\bf Source Cart} A cart intended only for use in the DCR with a surface suitable for working on sources and a recessed section that can hold the cleaning module. Should have drawers suitable for the storage of tools required for exchanging sources; the necessary tools should stay with the cart. 
\item {\bf Source Connector} A device by which a source may be temporarily connected to an umbilical while also allowing fibre optic or electrical connections.
\item {\bf URM (Umbilical Retrieval Mechanism)} The device in which the umbilical is stored in preparation to deployment that also includes drive mechanisms for the central rope and umbilical to facilitate vertical movement of sources.
\item {\bf Source Tube} A solid length of tube connecting to the bottom of the URM with a bellows below that terminates on a gate valve that connects to a nipple  assembly on the UI when the URM is in a position ready for deployment. 
\item {\bf Umbilical} A length of Tygothane tubing containing an optical fiber and five electrical wires to which calibration sources are attached for deployment in the AV.
\item {\bf Boil Off Nitrogen} Nitrogen sourced from a cryogenic nitrogen dewar that is installed in the Utility drift and is transported to the DCR by way of copper conduit (high pressure) or from the international dewar (low pressure). The pressure of the high pressure nitrogen boil off for purging purposes must be regulated using the gas board in the DCR to less than 20 psi.
\end{itemize}

\subsection{Description of the Cleaning Module}

The cleaning module consists of a 3' tube with a 10" interior diameter. Inlet and outlet gas connections are supplied to allow nitrogen flushing of the module independent of the URM system. Two LAB sprayers near the top of the cleaning module provides a stream of LAB to actively rinse off the source and provide a stream that can fill the cleaning module. A drain at the bottom of the cleaning module allows the LAB to be sampled. The flange at the top of the source allows the cleaning module to the gate valve at the bottom of the URM source bellows. A cart with a jack is provided to carry and lift the cleaning module into position, and provide a clean working surface for source work.

\section{Procedures}
The following procedures must be done within the deck clean room and all necessary protocols must be followed. It is assumed that the operator is to perform all operations using latex gloves in addition to standard PPE. Dust counts should be monitored prior to operations involving opening the  URM source tube or exposing a source. If the dust count appears to be high, the DCR should be cleaned and the intended procedure should be suspended. All procedures listed here should be completed with the participation of a local calibration expert. 

\subsection{Connecting the URM to the UI}

\subsection{Disconnecting the URM from the UI}

\subsection{Using the Source Connector}

\subsection{Attaching a Source to the Umbilical}\label{ss:sourceattach}
The instructions below presuppose that the source is not connected to the umbilical. Note that the following only applies to the North side URM as the South side URM will have the Cherenkov source permanently attached to its umbilical. Sources of LAB, and boil off nitrogen should be procured before attempting the following procedure. 

\begin{}


\begin{enumerate}
\item Begin with the source in the cleaning module
\item Start the flushing the cleaning module with N$_{2}$ 
\item Move the URM back from the UI over the working area West of the UI
\item Position the cleaning module cart under the URM
\item Close the URM cover gas bag and start actively flushing the URM with N$_{2}$
\item Open the gate valve and lower the source connector until it can reach the source.
\item Connect the source to the umbilical via the source connector. 
\item Raise the source into the source tube
\item Connect the cleaning module to the source tube gate valve
\item Connect the LAB supply to the cleaning module
\item Start the LAB flow and lower the source through the LAB stream. The source should be  raised and lowered through the spray to remove surface contamination. The initial rinse should continue until the source (and the source connector) are completely immersed in LAB.
\item Monitor the O$_2$ levels in the URM to validate the cover gas viability. The URM radon levels should also be monitored directly using the RAD7. The cover gas should be considered viable after the measured radon levels are consistent with zero. Once the cover gas is deemed viable the N$_{2}$ flush can be terminated and the cover gas bags can be opened to the URM. 
\item Drain the LAB from the cleaning module and take a sample to be submitted for analysis.
\item Repeat spraying the source with LAB as above until the source is again immersed in LAB. Allow the source to sit in the LAB for some time (order of a week) periodically draining the LAB, submitting samples for analysis and spraying the source until the source is again immersed. Record the status of the LAB cleanliness for each sample preparation (need to define a quantifiable cleanliness criteria)
\item Once the LAB cleanliness criteria is reached (TBD) the source should be raised into the source tube, the gate valve closed and the cleaning module disconnected from the gate valve. 
\item Move the URM back to the UI and reconnect the gate valve to the nipple assembly on the UI
\end{enumerate}
\subsection{Cleaning a Source Prior to Use}
Here, the minimum guidelines for cleaning a source should be set out. Each group should have their own specific cleaning instructions which should be compiled with this document. However, prior to use each source should
\begin{enumerate}
\item be inspected by a designated person to look for signs of corrosion or damage to the source
\item be wiped down using ultra pure water to initially pick up surface contaminants
\item be wiped down with an alconox mixture to futher remove surface contaminants
\item be wiped down with ultra pure water again to remove the alcanox
\end{enumerate}

\subsection{Detaching a Source from the Umbilical}
Required for this procedure is a storage module for the source. This module should be connected to a pump/purge system to keep the source in a nitrogen environment between uses.
\begin{enumerate}
\item Ensure that the source storage module is available for use on the source cart.
\item Disconnect the URM from the UI.
\item Move the URM back to the working area and place the source cart in the working area
\item Disconnect the URM cover gas bag and flush nitrogen through the URM.
\item Open the source tube gate valve over the working surface of the source cart.
\item Lower the source to the working surface of the source cart.
\item Disconnect the source from the umbilical at the source connector.
\item Place the source in the storage module. Close the storage module and pump nitrogen into the storage module.
\item If a new source is not to be immediately connected to the umbilical, retract the source connector back into the URM source tube and close the gate valve. Continue the URM purge until O$_{2}$ monitor and RAD7 values are below established thresholds before the cover gas bags are reconnected. 
\item If a new source is to be reconnected refer to Section \ref{ss:sourceattach}.
\end{enumerate}

\subsection{Procedure for the operation of LAB set up}
\begin{enumerate}
\item Make connections  for the  airline fitting at the specified port on the system for starting the pump.
\item Check the connection for the inlet and recirculation line at the LAB storage tank(Ensure that the tank is filled to the minimum level required for the cleaning).
\item Set the bottom MV1  toward the injection line and the MV4  toward  line to the vessel.
\item Start the pump by providing the pressure from the Air line ( Maximum upto 125 psi,pressure can be controlled by the pressure control valve) and allow the LAB flow in the system.
\item Wait until an adequate amount of fluid is pumped in the system. 
\item Once this is achieved, stop the pump and set the MV1 toward the recirculation  line.
\item Turn on the pump again and run the operation
\item Place the source in the storage module. Close the storage module and pump nitrogen into the storage module.
\item If a new source is not to be immediately connected to the umbilical, retract the source connector back into the URM source tube and close the gate valve. Continue the URM purge until O$_{2}$ monitor and RAD7 values are below established thresholds before the cover gas bags are reconnected. 
\item If a new source is to be reconnected refer to Section \ref{ss:sourceattach}.
\end{enumerate}


